% 부록. 결정기록 색인 · 재현 절차 · 선행연구 재현 상세

\section{결정기록}

설계를 정할 때마다 왜 그렇게 정했는지를 한 편씩 남겼다. 근거와 함께
\textbf{버린 대안}도 적는다. 본문에서 ``0014''와 같이 가리키는 것이 이 번호이며,
전문은 저장소의 \texttt{docs/decisions/}에 있다.

% 자동 생성 — scripts/build_tables.py. 손으로 고치지 않는다.
% 출처: docs/decisions/*.md
\begin{table}[htbp]
  \centering
  \caption{결정기록. 설계를 정할 때마다 왜 그렇게 정했는지를 한 편씩 남겼고 근거와 함께 \textbf{버린 대안}도 적었다.}
  \label{tab:decisions}
  \begin{tabular}{llp{9.5cm}}
    \toprule
    \textbf{번호} & \textbf{날짜} & \textbf{제목} \\
    \midrule
    0001 & 09-03 & 주 결과변수를 가격에서 거래량으로 변경 \\
    0002 & 09-03 & 분석 구간을 2025-10-19 이전으로 한정 \\
    0003 & 09-03 & 단지(aptSeq) 고정효과를 기본 사양으로 채택 \\
    0004 & 09-03 & 취소거래 제외 및 신고지연 절단 \\
    0005 & 09-03 & 처치 역전 구간에 de Chaisemartin \& D'Haultfœuille 추정량 사용 \\
    0006 & 09-03 & 모든 결과를 반사실 대비로 명시하여 서술한다 \\
    0007 & 09-07 & 문서를 동결·살아있는·최종으로 나누고, 숫자에 조건을 붙인다 \\
    0008 & 09-07 & 지정 유지 단지를 단지명이 아니라 지번으로 확정한다 \\
    0009 & 09-07 & 통제군 오염 점검 결과와 대응 \\
    0010 & 09-07 & 2020년 지정 범위를 법정동 단위 처치로 확정한다 \\
    0011 & 09-08 & 통제군을 어떻게 고르는가 \\
    0012 & 09-08 & 처치군을 무엇으로 하고, 무엇을 확장으로 미루는가 \\
    0013 & 09-11 & 사건 시점을 효력일에 맞추고, 거래 없는 달을 0으로 채운다 \\
    0014 & 09-11 & 거래량의 사전 추세 위배를 어떻게 다루는가 \\
    0015 & 09-11 & 대출 규제를 어떻게 다루는가 — 무엇이 상쇄되고 무엇이 남는가 \\
    0016 & 09-11 & 위약 검정 — 효과는 있으나 크기는 구간으로 말한다 \\
    0017 & 09-11 & 합성대조군 — 평균 회귀 의심에 답한다 \\
    0018 & 09-11 & 반복매매지수를 어떻게 짓는가 \\
    0019 & 09-18 & 클러스터가 적을 때의 추론 — 관행 p 를 버린다 \\
    0020 & 09-18 & 오염 동을 빼도 직후 효과는 그대로이고, 24개월 효과는 커진다 \\
    0021 & 09-18 & 15억에는 계단이 없다 — 가격 차이는 면적 면제와 겹친다 \\
    0022 & 09-18 & 이질성은 대출이 아니라 허가 면적 기준에서 왔다 \\
    0023 & 09-23 & 2025년 구간은 해제가 아니라 재지정으로 본다 \\
    0024 & 09-23 & 파급은 같은 동 안에만 있었다 — '인근 급등'은 기준점이 만든 말이다 \\
    0025 & 09-23 & 가격 효과는 −4 \textasciitilde{} −7\% 로 모인다 — 다만 합성대조군으로는 확인되지 않는다 \\
    0026 & 09-30 & 사전 추세를 감안하면 직후 거래량만 남는다 \\
    0027 & 09-30 & 음의 가중치는 작다 — 부과는 뚜렷하고 해제는 가릴 수 없다 \\
    0028 & 09-30 & 192개 사양 전부 음수다 — 갈린 것은 통제군이 아니라 날짜 세는 법이었다 \\
    0029 & 10-01 & 폐지 뒤 고가 거래는 돌아왔지만 그 규제 덕으로 돌릴 수 없다 \\
    0030 & 10-01 & 취소거래는 정책 반응이 아니다 — 그리고 0004 의 한계는 생각보다 작다 \\
    0031 & 10-08 & 강남·송파 밖에서도 같다 — B그룹이 A그룹과 다르다고 할 수 없다 \\
    \bottomrule
  \end{tabular}

  \begin{minipage}{\linewidth}\footnotesize
  \textit{주.} 날짜는 \textbf{작성일}이다 --- 뒤에 보강·정정된 기록이 있다. \textbf{이 표는 \texttt{docs/decisions/}의 파일들에서 생성된다.} 손으로 적었을 때는 색인의 날짜 셋이 정작 그 기록과 달랐다.
  \end{minipage}
\end{table}


\textbf{결정이 바뀐 것도 지우지 않고 그 자리에 정정으로 남겼다.} 0001의 근거 2는
2026-09-11에, 0004의 수치와 한계는 2026-10-01과 10-15에, 0015의 이질성 결론은
0021과 0022에서 바뀌었다. 바뀐 과정 자체가 이 연구의 기록이다.

\section{자료 출처와 이용 조건}

이 연구의 자료는 전부 공개 자료이며, 저장소를 공개하기 때문에 \textbf{무엇을
넣었고 무엇을 뺐는지}를 밝혀 둔다.

\begin{itemize}[leftmargin=2em, itemsep=3pt]
  \item \textbf{아파트 매매 실거래 상세 자료} --- 국토교통부, 공공데이터포털
        (\texttt{data.go.kr}) 오픈 API \texttt{RTMSDataSvcAptTradeDev}.
        활용신청 승인 뒤 이용했다. \textbf{인증키는 저장소에 넣지 않는다} ---
        \texttt{.env}(git 제외)에서 읽고 저장소에는 \texttt{.env.example}만 둔다.
        내려받은 원자료 902\,MB도 저장소 밖이다.
  \item \textbf{토지거래허가구역 지정·해제 공고} --- 서울특별시 \emph{서울시보}와
        전자관보에 실린 공고 원문. 누구나 열람할 수 있는 공보이며, 이 연구가 쓴
        공고 번호는 본문과 \texttt{data/policy/ltz\_events.csv}에 적었다.
        서울시보 원본(2.9\,GB)은 용량 때문에 저장소 밖에 두고 \textbf{해당
        공고만 발췌한 PDF}를 저장소에 넣었다.
  \item \textbf{행정구역 경계} --- 법정동 폴리곤. 맞닿은 동 판정(7장)에만 쓴다.
  \item \textbf{선행연구 논문} --- \textbf{저장소에 넣지 않았다.} 학술지 논문은
        저작권이 있으므로 개인적으로 내려받아 읽었을 뿐이며, 재현 코드
        (\texttt{scripts/replicate/})와 \textbf{논문이 보고한 수치만} 저장소에
        있다. 재현에는 원문 PDF가 필요하지 않다.
\end{itemize}

\textbf{개인정보는 다루지 않는다.} 실거래 자료는 단지·면적·층·계약일·금액만
담고 거래 당사자를 식별하는 정보가 없다. 가장 작은 분석 단위가
\textbf{(단지, 전용면적, 층)}이며 이는 특정 호실이 아니다(3.4절).

\section{재현 절차}

\subsection{전체 실행}

\begin{verbatim}
pip install -r requirements.txt
cp .env.example .env          # 공공데이터포털 Decoding 키 입력
python run_all.py --collect   # 수집부터 전부 (수십 분)
python run_all.py             # 수집 생략, 정제부터 문서까지 전부 (약 21분)
python run_all.py --list      # 전체 단계 목록
\end{verbatim}

\subsection{클린 환경에서 재현할 때}

새 폴더에 내려받아 새 가상환경에서 돌려 \texttt{output/results/*.json} 25개가
바이트 단위로 같은 것을 확인했다(2026-10-15). 그때 걸린 것 셋을 적어 둔다.

\textbf{첫째, \texttt{data/raw}는 git 밖이다}(902\,MB). \texttt{--collect}로 다시
받으면 \textbf{같은 결과가 나오지 않는다} --- 실거래 자료가 소급 갱신되기
때문이다. 재현하려면 \textbf{같은 스냅샷을 복사}해 와야 한다. 현재 스냅샷은
\textbf{2026-09-03 수집}이다.

\textbf{둘째, 경로가 길면 설치가 깨진다.} 윈도 기본 한계가 260자인데
\texttt{statsmodels}의 시험 자료 경로가 깊어 넘어간다. 137자짜리 폴더에서 266자가
되어 \texttt{statsmodels}만 빠진 채 설치가 끝났다. 짧은 경로에 두거나 긴 경로
지원을 켠다.

\textbf{셋째, 소요 시간.} \texttt{pip install} 약 12분, \texttt{run\_all.py} 약
21분.

\subsection{최종 보고서 빌드}

\begin{verbatim}
winget install --id MiKTeX.MiKTeX --scope user
.\paper\build.ps1
\end{verbatim}

\textbf{그림은 하나도 복사하지 않는다} --- \texttt{output/figures}를 그대로
참조하므로 분석을 다시 돌리면 보고서의 그림이 따라 바뀐다. 표는 전부가 아니라
\textbf{수치가 결과 JSON에 있는 \pipeTables 개}를
\texttt{scripts/build\_tables.py}가 생성한다
(\texttt{paper/tables/*.tex}). 나머지 표는 설계 비교·제도 요약처럼 분석 산출물이
아닌 것을 담으므로 본문에 직접 적었다. \textbf{손으로 적은 표에 분석 수치를
넣지 않는 것}이 경계선이다.

\section{선행연구 재현 상세}

\texttt{scripts/replicate/}에 전상현 외(2025)와 이한웅·이춘원(2025)의 사양을
복제한 코드가 있다. 본문 2.3절의 요약을 뒷받침하는 수치는 다음과 같다.

\begin{table}[htbp]
  \centering
  \caption{전상현 외(2025) 재현 --- 빈 셀 처리에 따라 교호항 부호가 바뀐다}
  \label{tab:replicate-jeon}
  \begin{tabular}{lrrrrr}
    \toprule
    \textbf{사양} & \textbf{Zone} & \textbf{Time} & \textbf{Zone\#Time} &
    \textbf{거래건수} & \textbf{$R^2$} \\
    \midrule
    논문 보고값 & $+0.3263$ & $+0.2724$ & $+0.2928$ & $+0.0116$ & 0.0757 \\
    빈 셀 0 코딩 ($N=675$) & $+0.3464$ & $+0.2334$ & $+0.3549$ & $+0.0138$ &
      0.0804 \\
    \textbf{빈 셀 제외} ($N=644$) & $+0.3038$ & $+0.1964$ &
      $\mathbf{-0.0448}$ & $+0.0023$ & 0.2496 \\
    \bottomrule
  \end{tabular}
\end{table}

빈 셀 31개는 전부 통제동(삼전·석촌·압구정·율현)이다. 거래건수를 통제에서 빼면
교호항이 $-0.1121$($p=0.002$), 동·월 고정효과까지 넣으면 $-0.1215$가 된다.

\begin{table}[htbp]
  \centering
  \caption{이한웅·이춘원(2025) --- 논문 표의 수치로 차이의 차이를 계산하면}
  \label{tab:replicate-lee}
  \begin{tabular}{lrrr}
    \toprule
    \textbf{창} & \textbf{잠실} & \textbf{비잠실} & \textbf{차이의 차이} \\
    \midrule
    $\pm$3개월 & $+16.5\%$ & $+4.3\%$ & $+12.2$\,\%p \\
    $\pm$6개월 & $+9.6\%$ & $+12.5\%$ & $\mathbf{-2.9}$\,\%p \\
    $\pm$12개월 & $+19.0\%$ & $+22.9\%$ & $\mathbf{-3.9}$\,\%p \\
    \bottomrule
  \end{tabular}
\end{table}

논문이 보고한 $t$-검정은 각 집단 \textbf{내부}의 사전-사후 비교이며 이중차분
계수가 아니다. 잠실동 가격이 유의하게 올랐다는 것은 맞지만 통제동도 함께, 그리고
$\pm$6개월부터는 더 많이 올랐다.
